\documentclass[12pt, a4paper]{article}
\usepackage[utf8]{inputenc}
\usepackage{amsmath, amssymb, amsthm, amsfonts,amsthm,tikz-cd}
\usepackage{marvosym}
\usepackage{mathrsfs}
\usepackage{CJKutf8}
\usepackage{bm}
\usepackage[margin=1in]{geometry}
\newcommand{\id}{\operatorname{id}}
\newcommand{\qz}{\mathbb{Q}/\mathbb{Z}}
\newcommand{\z}{\mathbb{Z}}
\renewcommand{\hom}{\operatorname{Hom}}
\theoremstyle{remark}
\newtheorem*{remark}{Remark}
\newenvironment{lemma}[1]{
    \noindent \textbf{Lemma #1.} \itshape
}{
    \vspace{1em}
}

% 重新定义 solution 环境，保留标识并在末尾留出 8cm 空白
\newenvironment{solution}
  {\paragraph{\textit{Solution:}}\par\medskip}
  {\hfill$\blacksquare$\vspace{0.1cm}}

\title{Abstract Algebra III: Assignment - Week 12}
\author{
    \begin{CJK*}{UTF8}{gbsn}
    钟皓宇 (12533556)
    \end{CJK*}
}
\date{\today}

\begin{document}

\maketitle
\section*{Problem 1}
\begin{enumerate}
    \item[(i)] Let $(X, \rho)$ be a transitive \textbf{Set}-representation of $G$. For $x \in X$, let $G_x$ denote the fixator of $x$ in $G$. Then the map
    $$
    G/G_x \to X, \quad gG_x \mapsto \rho(g)(x)
    $$
    defines an isomorphism $(G/G_x, \rho_{G_x}) \xrightarrow{\sim} (X, \rho)$.
    \item[(ii)] For subgroups $H$ and $H'$ of $G$, the representations $(G/H, \rho_H)$ and $(G/H', \rho_{H'})$ are isomorphic if and only if $H$ and $H'$ are conjugate in $G$.
\end{enumerate}

\begin{solution}
    (i) Let $f:G/G_x\to X$ denote the map we discuss. First we show that \(f\) is well-defined. Suppose
\[
gG_x=hG_x.
\]
Then \(h^{-1}g\in G_x\), so
\[
\rho_{h^{-1}g}(x)=x.
\]
Hence
\[
\rho_g(x)=\rho_h(\rho_{h^{-1}g}(x))=\rho_h(x).
\]
Therefore \(f(gG_x)=f(hG_x)\), so \(f\) is well-defined.
    Next, We show that $f$ is a bijection of sets. Let $y$ be an arbitrary element of $X$ that distinct to $x$. 
    Since $G$ act transitively on $X$ then there exists $g\in G$ such that $\rho_g(x)=y$. Therefore, we have 
    \begin{align*}
        f(gG_x) = \rho_g(x) = y.
    \end{align*}
    It implies that $f$ is surjective. Let $g,h\in G$ such that $f(gG_x) = f(hG_x)$. It implies that $\rho_g(x)=\rho_h(x)$ and thus 
    \begin{align*}
        x = \rho_{h^{-1}}\circ\rho_g(x) = \rho_{h^{-1}\cdot g}(x).
    \end{align*} 
    Then we have $h^{-1}\cdot g\in G_x$. Equivalently, we have $gG_x = hG_x$. Therefore, $f$ is injective. 
    For arbitrary $g\in G$, we then have 
    \begin{align*}
        f(g\cdot hG_x) = f(gh G_x) = \rho_{gh}(x) = \rho_g\circ\rho_h(x) = \rho_g \cdot f(hG_x).
    \end{align*}
    Therefore, $f$ is isomorphism of representation.

    (ii) Suppose first that
\[
f:G/H\to G/H'
\]
is an isomorphism of \(G\)-sets. Write
\[
f(H)=gH'
\]
for some \(g\in G\). Since \(f\) is \(G\)-equivariant, it preserves stabilizers. Hence
\[
H=\operatorname{Stab}_G(H)
=\operatorname{Stab}_G(f(H))
=\operatorname{Stab}_G(gH')
=gH'g^{-1}.
\]
Thus \(H\) and \(H'\) are conjugate.

Conversely, suppose
\[
H=gH'g^{-1}
\]
for some \(g\in G\). Define
\[
\Phi:G/H\to G/H',
\qquad
\Phi(aH)=agH'.
\]
This is well-defined: if \(aH=bH\), then \(a=bh\) for some \(h\in H\). Since
\[
g^{-1}hg\in H',
\]
we get
\[
agH'=bhgH'=bgH'.
\]
Thus \(\Phi(aH)=\Phi(bH)\).

The map \(\Phi\) is bijective, with inverse
\[
\Psi:G/H'\to G/H,
\qquad
\Psi(aH')=ag^{-1}H.
\]
Moreover, for \(k\in G\),
\[
\Phi(k\cdot aH)
=
\Phi(kaH)
=
kagH'
=
k\cdot \Phi(aH).
\]
Therefore \(\Phi\) is \(G\)-equivariant. Hence
\[
(G/H,\rho_H)\cong (G/H',\rho_{H'})
\]
if and only if \(H\) and \(H'\) are conjugate in \(G\).
\end{solution}



\section*{Problem 2}
Assume that $G$ is the cyclic (additive) group $\mathbb{Z}/p\mathbb{Z}$ where $p$ is a prime number, and let $k$ be a field of characteristic $p$. 
Consider the morphism
$$
G \to \mathrm{GL}_2(k), \quad \lambda \mapsto \begin{pmatrix} 1 & \lambda \\ 0 & 1 \end{pmatrix}.
$$
Prove that the line generated by the first standard basis vector of $k^2$ is $G$-stable but has no $G$-stable complement.
\begin{solution}
    Let
\[
e_1=\begin{pmatrix}1\\0\end{pmatrix},
\qquad
e_2=\begin{pmatrix}0\\1\end{pmatrix}.
\]
Since \(\operatorname{char}(k)=p\), the map
\[
\rho:\mathbb Z/p\mathbb Z\to \mathrm{GL}_2(k),
\qquad
\lambda\mapsto
\begin{pmatrix}
1&\lambda\\
0&1
\end{pmatrix}
\]
is well-defined and is a group homomorphism.

Let
\[
W=ke_1.
\]
Then for every \(\lambda\in \mathbb Z/p\mathbb Z\),
\[
\rho(\lambda)e_1=e_1,
\]
so \(W\) is \(G\)-stable.

We claim that \(W\) has no \(G\)-stable complement. Suppose, for contradiction, that
\[
k^2=W\oplus U
\]
with \(U\) is \(G\)-stable. Then \(U\) is one-dimensional and \(U\neq W\), so
\[
U=k(ae_1+e_2)
\]
for some \(a\in k\). Since \(U\) is \(G\)-stable, it is stable under the action of \(1\in \mathbb Z/p\mathbb Z\). Thus
\[
\rho(1)(ae_1+e_2)=(a+1)e_1+e_2
\]
must lie in \(U\). Hence there exists \(c\in k\) such that
\[
(a+1)e_1+e_2=c(ae_1+e_2).
\]
Comparing coefficients gives \(c=1\), and then
\[
a+1=a,
\]
a contradiction. Therefore \(W\) has no \(G\)-stable complement.
\end{solution}



\section*{Problem 3}
Let $(M, \rho)$ and $(N, \sigma)$ be two $k$-representations of $G$. Fix basis $(e_i)_{1 \le i \le m}$ and $(f_j)_{1 \le j \le n}$ of $M$ and $N$ respectively. Then $(e_i \otimes f_j)_{1 \le i \le m, 1 \le j \le n}$ is a basis of $M \otimes_k N$. Denote by $(m, \rho)$, $(n, \sigma)$, $(mn, \rho \otimes \sigma)$ be the matrix representations associated with $(M, \rho)$, $(N, \sigma)$, $(M \otimes_k N, \rho \otimes \sigma)$ respectively.
Determine the relation between $(\rho \otimes \sigma)(g)$ and $\rho(g)$, $\sigma(g)$.
\begin{solution}
    Suppose
\[
\rho_g(e_i)=\sum_{k=1}^m a_{ki}e_k,
\qquad
\sigma_g(f_j)=\sum_{\ell=1}^n b_{\ell j}f_\ell .
\]
Then
\[
(\rho\otimes\sigma)_g(e_i\otimes f_j)
=
\rho_g(e_i)\otimes\sigma_g(f_j)
=
\sum_{k=1}^m\sum_{\ell=1}^n
a_{ki}b_{\ell j}(e_k\otimes f_\ell).
\]
With respect to the ordered basis
\[
(e_1\otimes f_1,\ldots,e_1\otimes f_n,
e_2\otimes f_1,\ldots,e_m\otimes f_n),
\]
the column corresponding to \(e_i\otimes f_j\) is
\[
(a_{1i}b_{1j},\ldots,a_{1i}b_{nj},
a_{2i}b_{1j},\ldots,a_{2i}b_{nj},
\ldots,
a_{mi}b_{1j},\ldots,a_{mi}b_{nj})^T.
\]
Therefore
\[
[(\rho\otimes\sigma)_g]=[\rho_g]\otimes[\sigma_g],
\]
where the tensor product on the right is the Kronecker product of matrices.
\end{solution}



\section*{Problem 4}
Let $(M, \rho)$ be a $k$-representations of $G$. Fix basis $(e_i)_{1 \le i \le m}$ of $M$. Let $(M^*, \rho^*)$ be the contragredient representation of $(M, \rho)$. Denote by $(m, \rho)$, $(m, \rho^*)$ be the matrix representations associated with $(M, \rho)$, $(M^*, \rho^*)$.
Determine the relation between $\rho(g)$ and $\rho^*(g)$.

\begin{solution}
Let \(e_1,\dots,e_n\) be a basis of \(V\), with dual basis
\(e_1^*,\dots,e_n^*\). Suppose
\[
\rho_g(e_i)=\sum_{k=1}^n a_{ki}e_k.
\]
Thus
\[
[\rho_g]=(a_{ki}).
\]

The contragredient representation \(\rho^*\) is defined by
\[
(\rho_g^*\varphi)(v)=\varphi(\rho_{g^{-1}}v).
\]
Write
\[
\rho_{g^{-1}}(e_j)=\sum_{k=1}^n b_{kj}e_k.
\]
Then \([\rho_{g^{-1}}]=(b_{kj})=[\rho_g]^{-1}\). For each \(i,j\),
\[
(\rho_g^*e_i^*)(e_j)
=
e_i^*(\rho_{g^{-1}}e_j)
=
e_i^*\left(\sum_{k=1}^n b_{kj}e_k\right)
=
b_{ij}.
\]
Hence the coefficient of \(e_j^*\) in \(\rho_g^*(e_i^*)\) is \(b_{ij}\). Therefore
\[
[\rho_g^*]=(b_{ij})= [\rho_{g^{-1}}]^T=([\rho_g]^{-1})^T.
\]
Thus the matrix of the contragredient representation is the inverse transpose
of the original matrix.
\end{solution}


\section*{Problem 5}
\begin{enumerate}
    \item[(i)] Let $G$ be a finite group acting on a finite set $\Omega$. Let $\mathrm{Fix}^G(\Omega)$ denote the set of fixed points of $\Omega$ under $G$, and let $\Omega^\#$ denote the set of orbits of $G$ on $\Omega \setminus \mathrm{Fix}^G(\Omega)$. Prove that
    $$
    |\Omega| = |\mathrm{Fix}^G(\Omega)| + \sum_{\omega \in \Omega^\#} |\omega|.
    $$
    From now on, we assume that $p$ is a prime number and that $G$ is a nontrivial $p$-group.
    \item[(ii)] Assume now that $|\Omega|$ is a power of $p$ different from 1. Prove that $|\mathrm{Fix}^G(\Omega)|$ is divisible by $p$.
    \item[(iii)] Let $k$ be a (commutative) field of characteristic $p$. Let $M$ be a $kG$-module. Prove that $\mathrm{Fix}^G(M) \neq 0$. (Hint. (a) Let $(e_i)_{i=1,\dots,r}$ be a $k$-basis of $M$. Prove that the abelian group generated by $(g e_i)_{i=1,\dots,r \ (g \in G)}$ is an $\mathbb{F}_p G$-module. (b) Apply question (2) above to conclude.)
    \item[(iv)] Prove that, up to isomorphism, the trivial representation of $G$ is the unique irreducible $kG$-module.
\end{enumerate}
\begin{solution}
    (i) The set \(\Omega\) decomposes as a disjoint union
\[
\Omega=\operatorname{Fix}^G(\Omega)\sqcup
\bigl(\Omega\setminus \operatorname{Fix}^G(\Omega)\bigr).
\]
Moreover, \(\Omega\setminus \operatorname{Fix}^G(\Omega)\) is \(G\)-stable, hence it is a disjoint union of its \(G\)-orbits:
\[
\Omega\setminus \operatorname{Fix}^G(\Omega)
=
\bigsqcup_{\omega\in \Omega^\#}\omega .
\]
Therefore
\[
\Omega
=
\operatorname{Fix}^G(\Omega)
\sqcup
\bigsqcup_{\omega\in \Omega^\#}\omega .
\]
Taking cardinalities gives
\[
|\Omega|
=
|\operatorname{Fix}^G(\Omega)|
+
\sum_{\omega\in \Omega^\#}|\omega|.
\]

(ii) We just prove that for any $\omega\in \Omega^\#$ we have $|\omega|$ is divisible by $p$. Consider $\omega$ as transitive $G$-set, according to (1), we choose $x\in \omega$ then it induces the bijection
\begin{align*}
    G/G_x \xrightarrow{\sim} \omega.
\end{align*}
Since the only prime divisor of $|G|$ is $p$, we can suppsoe that $|\omega|= |G/G_x|=p^s.$ If $s=0$, it turns out that $x\in \operatorname{Fix}^G(\Omega)$, which contradicts to our assumption.
Therefore, $|\Omega|$ and each $|\omega|$ is divisible by $p$ implies that 
\begin{align*}
|\operatorname{Fix}^G(\Omega)| =  |\Omega| - \sum_{\omega\in \Omega^\#}|\omega|.
\end{align*}
is also divisible by $p$.

(iii) Assume \(M\neq 0\), and let \(e_1,\dots,e_r\) be a \(k\)-basis of \(M\).
Set
\[
M'=\langle ge_i\mid g\in G,\ 1\leq i\leq r\rangle_{\mathbb F_p}\subset M.
\]
Since \(\operatorname{char}(k)=p\), this is a finite \(\mathbb F_p\)-subspace of
\(M\). It is \(G\)-stable, because for \(h\in G\),
\[
h(ge_i)=(hg)e_i\in M'.
\]
Thus \(G\) acts on the finite set \(M'\).

Since \(M'\neq 0\), we have \(|M'|=p^n\) for some \(n\geq 1\). By part (ii),
\[
|\operatorname{Fix}^G(M')|
\]
is divisible by \(p\). But \(0\in \operatorname{Fix}^G(M')\), so the fixed
point set cannot consist only of \(0\). Hence there exists
\[
0\neq m\in \operatorname{Fix}^G(M').
\]
Since \(M'\subset M\), we get
\[
0\neq m\in \operatorname{Fix}^G(M).
\]
Therefore
\[
\operatorname{Fix}^G(M)\neq 0.
\]

(iv)Let \(M\) be an irreducible \(kG\)-module. By part (iii),
\[
\operatorname{Fix}^G(M)\neq 0.
\]
Choose
\[
0\neq m\in \operatorname{Fix}^G(M).
\]
Then \(gm=m\) for all \(g\in G\). Hence the one-dimensional subspace
\[
km\subset M
\]
is \(G\)-stable. Since \(M\) is irreducible, we must have
\[
M=km.
\]
Thus \(M\) is one-dimensional and \(G\) acts trivially on \(M\). Therefore
\[
M\cong k
\]
as \(kG\)-modules, where \(k\) is the trivial representation.

Conversely, the trivial representation \(k\) is one-dimensional, hence
irreducible. Therefore, up to isomorphism, the trivial representation is the
unique irreducible \(kG\)-module.


\end{solution}
\end{document}